%% Elevark fra mattelab.no – laget automatisk av bygg.py. Rediger gjerne til eget bruk.
%% Kompiler med lualatex.
\documentclass[11pt,a4paper]{article}
\usepackage[norsk]{babel}
\usepackage[a4paper,margin=2.2cm,top=2cm,bottom=2.4cm]{geometry}
\usepackage{fontspec}
\usepackage{unicode-math}
\setmainfont{TeX Gyre Pagella}
\setsansfont{TeX Gyre Heros}[Scale=MatchLowercase]
\setmathfont{TeX Gyre Pagella Math}
\usepackage{xcolor}
\definecolor{blekk}{HTML}{16213A}
\definecolor{dempet}{HTML}{566077}
\definecolor{rute}{HTML}{CFC9B8}
\definecolor{m1}{HTML}{2F4AB0}
\definecolor{m2}{HTML}{99266D}
\usepackage[most]{tcolorbox}
\usepackage{qrcode}
\usepackage{needspace}
\usepackage{fancyhdr}
\usepackage[hidelinks]{hyperref}
\pdfvariable suppressoptionalinfo \numexpr 1+2+4+8+16+32+64+128+256+512\relax
\setlength{\parindent}{0pt}
\setlength{\parskip}{0.4em}
\pagestyle{fancy}
\fancyhf{}
\renewcommand{\headrulewidth}{0pt}
\fancyfoot[L]{\sffamily\footnotesize\color{dempet}mattelab.no · Omvendte funksjoner og deriverte · elevark uten fasit}
\fancyfoot[R]{\sffamily\footnotesize\color{dempet}Side \thepage}

\newcommand{\fase}[2]{\tcbox[on line, colframe=#2, colback=white, coltext=#2, boxrule=0.6pt, arc=2pt,
  boxsep=0pt, left=3pt, right=3pt, top=1.5pt, bottom=1.5pt, fontupper=\sffamily\bfseries\footnotesize]{#1}}
\newcommand{\steg}[4]{%
  \par\noindent{\color{rute}\rule{\linewidth}{0.4pt}}\par\nopagebreak
  \noindent\begin{minipage}[t]{2.1cm}\vspace{0pt}\fase{#1}{#2}\end{minipage}%
  \begin{minipage}[t]{\dimexpr\linewidth-2.1cm\relax}\vspace{0pt}#3\end{minipage}\par
  \vspace{#4}}
\newcommand{\del}[4]{%
  \par\noindent{\color{rute}\rule{\linewidth}{0.4pt}}\par\nopagebreak
  \noindent\begin{minipage}[t]{1.1cm}\vspace{0pt}\fase{#1}{#2}\end{minipage}%
  \begin{minipage}[t]{\dimexpr\linewidth-1.1cm\relax}\vspace{0pt}#3\end{minipage}\par
  \vspace{#4}}
\newcommand{\oppgave}[2]{\par\Needspace{12\baselineskip}\vspace{0.6em}%
  {\fontsize{14}{17}\selectfont\bfseries{\color{m1}#1}\enspace #2}\par\vspace{0.2em}}

\begin{document}
\begin{minipage}[t]{0.72\linewidth}\vspace{0pt}
{\sffamily\small\color{dempet}Elevark · R1 · mattelab.no}\par\vspace{0.3em}
{\fontsize{24}{28}\selectfont\bfseries Omvendte funksjoner og deriverte}\par\vspace{1em}
{\sffamily\small Navn: \rule{4.6cm}{0.4pt}\quad Klasse: \rule{1.5cm}{0.4pt}\quad Dato: \rule{1.9cm}{0.4pt}}
\end{minipage}\hfill
\begin{minipage}[t]{0.22\linewidth}\vspace{0pt}\raggedleft
\qrcode[height=2.3cm]{https://mattelab.no/utforskere/omvendte-funksjoner.html}\par\vspace{0.2em}
{\sffamily\scriptsize\color{dempet}Utforskeren}
\end{minipage}

\vspace{1em}
{\sffamily\small\color{dempet}Utforskeren ligger på \url{https://mattelab.no/utforskere/omvendte-funksjoner.html}}\par
Hver oppgave har tre steg. \textbf{Gjett} først, uten å bruke utforskeren. \textbf{Test} så i utforskeren. \textbf{Forklar} til slutt med egne ord hvorfor det blir slik.

\oppgave{1}{Den deriverte av $\ln x$}
\steg{Gjett}{m1}{Grafen til $e^x$ har stigningstall $e \approx 2{,}72$ i $P = (1, e)$. Hva tror du stigningstallet til $\ln x$ er i det speilede punktet $Q$?}{2.2cm}
\steg{Test}{m2}{Velg $e^x$ og $\ln x$ med $a = 1$. Prøv også $a = 0$ og $a = -1$. Sammenlign stigningstallet i $Q$ med $\frac{1}{x}$, der $x$ er førstekoordinaten til $Q$.}{2.6cm}
\steg{Forklar}{blekk}{Bruk speilingen til å forklare at $(\ln x)' = \frac{1}{x}$.}{3.4cm}

\oppgave{2}{Kvadratroten}
\steg{Gjett}{m1}{$f(x) = x^2$ har $f'(2) = 4$. Hva er stigningstallet til $\sqrt{x}$ i $x = 4$?}{2.2cm}
\steg{Test}{m2}{Velg $x^2$ og $\sqrt{x}$ med $a = 2$. Dra så $a$ ned mot $0$.}{2.6cm}
\steg{Forklar}{blekk}{Hvorfor er ikke $\sqrt{x}$ deriverbar i $x = 0$?}{3.4cm}

\oppgave{3}{Tredjeroten}
\steg{Gjett}{m1}{$f(x) = x^3$ har en omvendt funksjon for alle $x$. Har $\sqrt[3]{x}$ en tangent i alle punkter? Hva skjer i origo?}{2.2cm}
\steg{Test}{m2}{Velg $x^3$ og $\sqrt[3]{x}$ med $a = 1$. Prøv $a = 2$ og $a = 0$.}{2.6cm}
\steg{Forklar}{blekk}{Finn stigningstallet til $\sqrt[3]{x}$ i $x = 8$ uten utforskeren.}{3.4cm}

\oppgave{4}{Lineære funksjoner}
\steg{Gjett}{m1}{Hva er den omvendte funksjonen til $f(x) = 2x + 1$? Hva er stigningstallet?}{2.2cm}
\steg{Test}{m2}{Velg $2x + 1$ med $a = 0$, og dra $a$ fram og tilbake.}{2.6cm}
\steg{Forklar}{blekk}{Vis at $g(f(x)) = x$. Hvor skjærer grafene hverandre?}{3.4cm}

\end{document}
